\documentclass[notitlepage,12pt]{jedm}

\usepackage[table]{xcolor}
\usepackage{url}
\usepackage{hyperref}
\usepackage{graphicx}
\usepackage{amsmath,amssymb,amsfonts}
\usepackage{array}
\usepackage{booktabs}
\usepackage{multirow}
\usepackage[hyphens]{xurl}
\hypersetup{
  colorlinks = true,
  urlcolor   = blue,
  linkcolor  = blue,
  citecolor  = blue
}

\begin{document}

\title{Measuring the Estimator, Not the Learner: A Confound in Simulated Evaluations of Adaptive Difficulty Selection}
\date{}

\author{{\large Kavya Aggarwal}\\Department of Computer Science and Engineering\\GITAM (Deemed to be University), Hyderabad, Telangana, India\\aggarwalkaavya23@gmail.com}

\maketitle

\begin{abstract}
Adaptive tutoring systems that select question difficulty are routinely evaluated against simulated learners, because running a classroom study for every design decision is impractical. A common construction fits Bayesian Knowledge Tracing (BKT) parameters per skill, generates responses from a latent knowledge state, and reports the BKT posterior as the learning outcome. We show that this construction cannot distinguish a policy that teaches more effectively from one that merely asks harder questions. The posterior is a function of the evidence it is shown, and difficulty determines that evidence: a correct response to a low-guess item is stronger evidence of mastery than a correct response to a high-guess item, whether or not any learning occurred. We make three contributions. First, we derive the mastery level at which each difficulty tier sits at a given target success rate, and show that for the guess and slip values in standard use these levels are $0.846$, $0.929$ and $1.000$ -- compressed into the top of the mastery scale, with the hardest tier optimal only in the limit of complete mastery. Over the range learners actually occupy the best action is therefore constant, and no adaptive policy can outperform the best fixed one. Second, we construct a generative model in which learning does not depend on difficulty, so that every policy produces an identical latent gain by construction, and measure what the posterior reports: routing toward harder items overstates the gain by up to $19\%$ of its true value at short horizons. Third, we show that the cost of mis-routing is asymmetric but that its direction depends on the assumed learning model: the worst case of serving items one tier too hard is a $89\%$ loss of attainable gain against $29\%$ for one tier too easy, making erring toward easier items the minimax choice under model uncertainty rather than a universally safe one. We conclude with a measurement protocol -- evaluate on a difficulty-neutral probe, and never let the controller and the outcome metric share an estimator -- and report that against a well-chosen fixed curriculum, adaptive routing requires roughly ten items per learner before it confers any advantage. All results are reproducible from released, seeded code.\\

{\parindent0pt
\textbf{Keywords:} Bayesian knowledge tracing, simulated learners, evaluation methodology, adaptive difficulty, measurement validity, contextual bandits
}
\end{abstract}

\section{Introduction}

Deciding what to ask a learner next is among the oldest problems in intelligent tutoring, and difficulty selection is one of its most studied instances. A tutor that can estimate what a learner knows should, in principle, be able to choose items that are neither trivial nor hopeless. Bayesian Knowledge Tracing \cite{corbett1995knowledge} supplies the estimate, and a bandit or similar policy supplies the choice.

Evaluating such a policy is harder than building it. A classroom study is expensive, slow, and cannot be run for every design variant, so the field relies heavily on simulated learners. A recent systematic review finds that simulated learners are widely used in educational technology and, more pointedly, that their validation is frequently absent or informal \cite{kaser2024simulated}. The present paper concerns a specific way in which a widely used simulation design can produce a result that looks like improved learning but is not.

The construction we examine is common and reasonable-looking. A latent binary knowledge state is generated per skill. Responses are emitted from that state through guess and slip parameters, with difficulty expressed by varying those parameters. BKT observes the responses and maintains a posterior. The policy under test selects difficulty. Finally -- and this is the step that matters -- the BKT posterior is read out and reported as the learner's mastery.

Our argument is that the final step silently changes what is being measured. The posterior is an inference about the learner given the evidence presented. Difficulty determines what evidence is presented. A correct answer on an item with a low guess parameter raises the posterior more than a correct answer on an item with a high guess parameter, because it is, correctly, stronger evidence. This is not a flaw in BKT; it is BKT working exactly as intended. It becomes a flaw only when the posterior is used as the outcome variable in a comparison between policies that differ in the difficulty they serve, because then part of the measured difference is inference rather than learning.

Whether this matters quantitatively depends on the generative model. If the simulation lets difficulty affect the latent learning rate, routing can genuinely change what is learned and the confound is partial. If the simulation fits a single learn rate per skill -- as is standard when parameters come from a fitting library \cite{bulut2023pybkt} -- then difficulty cannot affect learning at all, and the entire measured difference between a difficulty-adaptive and a difficulty-agnostic policy is attributable to the estimator. We show this case is not hypothetical and quantify its size.

\paragraph{Contributions.} We contribute: (i) an analytic characterisation of where each difficulty tier is pedagogically optimal under guess/slip encoding, showing that standard parameter values compress these points into the top of the mastery range (Section~\ref{sec:geometry}); (ii) a controlled demonstration that the reported gain separates by difficulty when the true gain cannot, with magnitudes (Section~\ref{sec:inflation}); (iii) a measurement of the asymmetric cost of mis-routing, showing that its direction is specification-dependent and that erring toward easier items is the minimax choice (Section~\ref{sec:asymmetry}); and (iv) a measurement protocol, together with the observation that adaptive routing needs a non-trivial horizon before it beats a well-chosen fixed curriculum (Sections~\ref{sec:horizon} and \ref{sec:protocol}).

\section{Background and Related Work}

\paragraph{Knowledge tracing and item difficulty.} BKT models a skill as a two-state hidden Markov process with parameters for prior knowledge, learning, guessing and slipping \cite{corbett1995knowledge}. In its standard form it has no item-difficulty parameter: all items exercising a skill are exchangeable. That this is a limitation is long established. KT-IDEM introduced per-item difficulty by giving each item its own guess and slip parameters \cite{pardos2011ktidem}, and the broader point that BKT parameters are not always identifiable from response data was made by Beck and Chang \cite{beck2007identifiability}, who showed that different parameter sets can fit the same data equally well while implying very different pedagogy. Our Section~\ref{sec:geometry} is a consequence of representing difficulty through guess and slip specifically, and is to our knowledge not stated elsewhere.

\paragraph{Difficulty selection policies.} Posterior sampling \cite{thompson1933likelihood} and its descendants are a natural fit for selecting among a small set of difficulty tiers, and systems combining knowledge tracing with bandit-style selection have been built and evaluated on simulated cohorts. Castleman et al.\ \cite{castleman2024hierarchical} present a hierarchical multi-armed bandit tutor that selects both concept and problem difficulty, evaluate it on simulated students whose responses are generated through BKT, and report that adding problem-difficulty adaptation improves mastery relative to a difficulty-agnostic variant. We return to this evaluation in Section~\ref{sec:inflation}, not to dispute its conclusion, which we are not in a position to re-derive, but because its stated construction is precisely the one our analysis concerns.

\paragraph{Difficulty and learning.} That difficulty has an interior optimum is a long-standing claim. The desirable-difficulties literature argues that conditions which slow acquisition can improve retention \cite{schmidt1992new}. More recently, Wilson et al.\ \cite{wilson2019eighty} derive an optimal training accuracy of approximately $85\%$ for a broad class of gradient-descent learners. We use that value as a pre-specified parameter where a target success rate is required, while noting explicitly that their demonstrations concern artificial and biologically plausible neural networks rather than human classroom learning, and that we treat it as a principled starting value rather than as evidence about students.

\section{Setting and Notation}
\label{sec:setting}

A learner has latent mastery $\theta \in [0,1]$ for a skill. Items belong to a difficulty tier $d$ with guess $g_d$ and slip $s_d$. The probability of a correct response is
\begin{equation}
P(\text{correct} \mid \theta, d) = \theta\,(1 - s_d) + (1 - \theta)\,g_d .
\label{eq:response}
\end{equation}
Throughout we use the tier values in common use and adopted by the system that motivated this work: $(g,s) = (0.30, 0.05)$ for easy, $(0.20, 0.10)$ for medium and $(0.10, 0.15)$ for hard.

The BKT posterior update on observing a response is the standard one: an evidence step conditioning on correctness through $g_d$ and $s_d$, followed by a transition step applying the learn rate. We write the estimator's belief as $\hat{\theta}$ to keep it typographically distinct from $\theta$. The distinction is the subject of this paper.

A policy maps the information available to the tutor -- in general some function of $\hat{\theta}$ -- to a tier. We compare policies by normalised learning gain, and in Sections~\ref{sec:inflation} onward we compute it twice: once from $\theta$, which a real evaluation cannot observe, and once from $\hat{\theta}$, which is what simulated evaluations typically report.

\section{Where Each Difficulty Tier Is Optimal}
\label{sec:geometry}

Suppose, following the interior-optimum view of difficulty, that learning is fastest when the probability of a correct response is close to some target $P^{*}$. Setting Equation~\ref{eq:response} equal to $P^{*}$ and solving for mastery gives the level at which tier $d$ sits exactly at the target:
\begin{equation}
\theta^{*}_d = \frac{P^{*} - g_d}{(1 - s_d) - g_d}.
\label{eq:thetastar}
\end{equation}

Two consequences follow immediately. First, the maximum attainable success probability on tier $d$ is $1 - s_d$, reached only at $\theta = 1$. A tier whose maximum lies at or below $P^{*}$ is therefore never strictly optimal for any attainable mastery. Second, $\theta^{*}_d$ is determined entirely by the tier parameters and the target.

Table~\ref{tab:geometry} evaluates Equation~\ref{eq:thetastar} at $P^{*} = 0.85$.

\begin{table}[ht]
\centering
\caption{Mastery at which each tier sits at the target success rate, for standard tier parameters and $P^{*}=0.85$.}
\label{tab:geometry}
\begin{tabular}{lcccc}
\toprule
Tier & $g_d$ & $s_d$ & Max.\ accuracy $(1-s_d)$ & $\theta^{*}_d$ \\
\midrule
Easy   & 0.30 & 0.05 & 0.950 & 0.846 \\
Medium & 0.20 & 0.10 & 0.900 & 0.929 \\
Hard   & 0.10 & 0.15 & 0.850 & 1.000 \\
\bottomrule
\end{tabular}
\end{table}

All three points lie in the top sixth of the mastery scale. The hard tier's maximum accuracy equals $P^{*}$ exactly, so it is optimal only in the limit $\theta \to 1$: a learner must have complete mastery before the hardest tier becomes the right thing to serve. The practical consequence is that over the mastery range a cohort actually occupies during instruction, the optimal action does not change. An adaptive policy has nothing to adapt to, and cannot be expected to beat the best fixed policy.

We stress what this result is and is not. It is not a claim that difficulty is unimportant, nor that adaptive selection is futile in general. It is a statement about a representation: encoding difficulty through guess and slip alone places the tiers' useful operating points close together and close to ceiling. Representations that give difficulty its own location parameter, as KT-IDEM does \cite{pardos2011ktidem}, are not subject to it. The result is also falsifiable and cheap to check: any system using tiered difficulty can evaluate Equation~\ref{eq:thetastar} on its own parameters before concluding that its routing policy underperformed for algorithmic reasons.

\section{The Estimator Separates When the Learner Does Not}
\label{sec:inflation}

We now isolate the confound. Consider a generative model in which the latent learning increment does not depend on which tier is served:
\begin{equation}
\theta_{t+1} = \theta_t + \gamma\,(1 - \theta_t),
\label{eq:indep}
\end{equation}
with $\gamma$ a learner-specific rate. This is not an artificial choice. It is exactly what obtains whenever BKT's learn parameter is fitted per skill rather than per difficulty, which is the default behaviour of standard fitting tools \cite{bulut2023pybkt} and the construction described by Castleman et al.\ \cite{castleman2024hierarchical}, whose appendix specifies a single learn parameter per skill with difficulty entering only through the response model.

Under Equation~\ref{eq:indep} the latent gain after a fixed number of items is identical for every policy, by construction. Any separation in reported gain is therefore entirely attributable to the estimator. Table~\ref{tab:inflation} reports both quantities for a cohort of 60 simulated learners across 200 trials, with mastery initialised in three ability bands.

\begin{table}[ht]
\centering
\caption{Latent and estimated gain under difficulty-independent learning (estimator learn rate $0.10$). The latent gain is identical for all tiers by construction; the estimated gain is not.}
\label{tab:inflation}
\begin{tabular}{cccccr}
\toprule
Items & Latent gain & Est.\ (easy) & Est.\ (medium) & Est.\ (hard) & Hard $-$ easy \\
\midrule
2  & 0.2562 & 0.1204 & 0.1340 & 0.1665 & $+0.0460$ \ (18.0\%) \\
3  & 0.3375 & 0.1989 & 0.2290 & 0.2630 & $+0.0641$ \ (19.0\%) \\
5  & 0.4456 & 0.3528 & 0.3854 & 0.4267 & $+0.0739$ \ (16.6\%) \\
8  & 0.5321 & 0.5044 & 0.5275 & 0.5583 & $+0.0540$ \ (10.1\%) \\
12 & 0.5843 & 0.5920 & 0.6042 & 0.6147 & $+0.0227$ \ (3.9\%) \\
\bottomrule
\end{tabular}
\end{table}

A policy that routes toward harder items reports up to $19\%$ more learning gain than one routing toward easier items, while producing exactly none. The mechanism is the evidence term: with $g_{\text{hard}} = 0.10$, a correct response is strong evidence of mastery, and the posterior moves accordingly. The effect is largest early and decays as the assumed learn rate drives the posterior toward ceiling, which has the uncomfortable implication that it is most severe precisely in the short-horizon regime where simulated comparisons are cheapest to run.

\paragraph{Scope of the claim.} We did not re-run the simulation of Castleman et al.; their repository publishes the tutoring algorithm and output trajectories but not the simulation harness, and the released trajectories record estimated mastery only, with no latent-state column. We therefore make the narrower claim that is actually supported: under the generative assumptions their appendix states -- a per-skill learn rate, difficulty entering through guess and slip, and mastery read from the BKT posterior -- a difference between difficulty-adaptive and difficulty-agnostic arms is not by itself evidence of improved learning. Their design also varies concept selection and advancement timing, which can alter exposure and is a separate pathway we do not quantify here. Our point is about what the metric can and cannot distinguish, not about the standing of any particular reported result.

\paragraph{A second-order effect.} When the same estimator both drives the controller and supplies the outcome metric, inflation does not merely flatter the score. If advancement to the next concept is triggered by estimated mastery, an inflated estimate advances learners sooner, on less actual knowledge. The metric and the control loop then fail in the same direction at the same time, and the system looks better precisely as it serves learners worse.

\section{The Cost of Mis-Routing Is Asymmetric, but Not Uniformly}
\label{sec:asymmetry}

A mastery estimator is usually judged by how close $\hat{\theta}$ is to $\theta$. For a system that acts on the estimate, what matters instead is what its errors cost, and the two directions of error need not cost the same. To measure this without entangling it with estimator quality, we compare policies that deliberately serve exactly one tier above, and one tier below, the tier that is optimal for the learner's true mastery at each step. Both make errors on every item; they differ only in direction.

\begin{table}[ht]
\centering
\caption{Learning lost by serving one tier above or below the optimum, relative to always serving the optimum, under three learning-model specifications.}
\label{tab:asymmetry}
\footnotesize
\begin{tabular}{llrrr}
\toprule
Horizon & Model & Optimal & Too easy & Too hard \\
\midrule
\multirow{3}{*}{6 items}
 & A: target, guess/slip & 0.1273 & $-0.0\%$  & $-74.9\%$ \\
 & B: target, IRT        & 0.3613 & $-0.9\%$  & $-89.0\%$ \\
 & C: difficulty matching        & 0.4626 & $-29.2\%$ & $-23.7\%$ \\
\midrule
\multirow{3}{*}{24 items}
 & A: target, guess/slip & 0.3136 & $-0.2\%$  & $-45.9\%$ \\
 & B: target, IRT        & 0.6076 & $-0.5\%$  & $-77.5\%$ \\
 & C: difficulty matching        & 0.6128 & $-14.7\%$ & $-1.5\%$ \\
\bottomrule
\end{tabular}
\end{table}

The asymmetry is large but its direction is not universal. Where learning is fastest at a target success rate (A and B), serving one tier too hard is severe -- costing between $46\%$ and $89\%$ of the attainable gain -- while serving one tier too easy costs almost nothing, because an over-easy item still yields practice whereas an over-hard one mostly yields failure. Where learning instead depends on matching difficulty to current mastery (C), the asymmetry weakens and reverses: over-easy items are then the more wasteful error.

A system designer does not know which specification is correct, so the relevant quantity is the worst case of each direction. Across the specifications and horizons tested, the worst case of serving too hard is a $89\%$ loss, against $29\%$ for serving too easy -- roughly three times greater. Erring toward easier items is therefore the minimax choice under model uncertainty, not because it is always better, but because its downside is bounded more tightly.

Two cautions. First, this is a statement about relative costs under stated generative assumptions, not an empirical claim about learner motivation or disengagement, which our simulation does not model. Second, the result should not be read as support for the desirable-difficulties position \cite{schmidt1992new} or against it: that literature concerns difficulty aiding long-term retention, which we do not measure, whereas this is about the immediate cost of mis-targeting.

The practical implication is narrower than a claim about estimator bias. An estimator for an acting system should be selected by the quality of the decisions it induces, and where the direction of the cost asymmetry is unknown, a conservative estimate is the safer default.

\section{How Long Before Adaptivity Pays}
\label{sec:horizon}

Adaptive policies are often compared against a random baseline, which is a low bar. Compared instead against the best fixed curriculum, adaptive routing is not uniformly better. Table~\ref{tab:horizon} reports the worst case across three learning-model specifications, scored against the oracle because the attainable gain differs between specifications.

\begin{table}[ht]
\centering
\caption{Worst-case attainment across three learning-model specifications, as a percentage of the oracle policy.}
\label{tab:horizon}
\begin{tabular}{cccc}
\toprule
Items per learner & Adaptive routing & Best fixed curriculum & Difference \\
\midrule
4  & 55.0 & 60.8 & $-5.8$ \\
6  & 60.4 & 61.2 & $-0.8$ \\
10 & 73.4 & 64.1 & $+9.3$ \\
16 & 76.3 & 67.9 & $+8.4$ \\
24 & 79.8 & 71.9 & $+7.9$ \\
50 & 84.2 & 79.1 & $+5.1$ \\
\bottomrule
\end{tabular}
\end{table}

Below roughly ten items per learner, a fixed curriculum is as good or better: the policy has not yet accumulated enough evidence to route well, and the cost of exploring is not recovered. This has a direct bearing on study design. A short simulated or classroom episode can report that adaptivity does not help, and be correct about that episode while saying nothing about the regime the system is intended for.

\section{A Measurement Protocol}
\label{sec:protocol}

The confound is straightforward to avoid, and the remedies are cheap.

\begin{enumerate}
\item \textbf{Do not report the posterior as the outcome.} Where a latent state exists in the simulation, measure it. It is already available and costs nothing to record.
\item \textbf{Where no latent state exists, measure on a difficulty-neutral probe.} Administer a fixed set of reference-difficulty items to every arm and score on those. Because the probe is identical across arms, its score is not confounded with the policy's choices.
\item \textbf{Do not let the controller and the metric share an estimator.} If advancement decisions and the reported outcome both derive from $\hat{\theta}$, errors in $\hat{\theta}$ move both in the same direction.
\item \textbf{State the difficulty--learning coupling explicitly.} If the generative model's learn rate does not depend on difficulty, then no difficulty policy can change learning, and this should be stated rather than left for the reader to infer from an appendix.
\item \textbf{Check Equation~\ref{eq:thetastar} before attributing failure to the algorithm.} If the tiers' optimal points are compressed into a region learners do not reach, the representation is the problem.
\end{enumerate}

\section{Limitations}

Our evidence is simulated. This is appropriate for the claims made -- that a measurement cannot distinguish two things is an analytical property, and the demonstration in Section~\ref{sec:inflation} is a controlled construction rather than an empirical finding about students -- but it does mean we make no claim about the magnitude of the confound in any deployed system with real learners.

The numeric results depend on the tier parameters in Section~\ref{sec:setting}. Those values are in common use, and Equation~\ref{eq:thetastar} is general, but systems using different parameters will obtain different $\theta^{*}_d$ and should compute their own.

Sections~\ref{sec:asymmetry} and \ref{sec:horizon} depend on a model in which difficulty affects learning, and the magnitude of both effects varies with the specification chosen. We report the worst case across three specifications rather than the most favourable, but the qualitative findings are specification-dependent in a way the Section~\ref{sec:geometry} result is not.

Finally, we have characterised a confound rather than established how often it is consequential in the published literature. Determining that would require re-running simulations that are frequently not released, which is itself a finding about reproducibility in this area.

\section{Conclusion}

Simulated learners make it affordable to evaluate instructional policies, but they also make it easy to measure the wrong thing. When the outcome variable is a knowledge-tracing posterior and the policy under test selects item difficulty, the comparison is confounded: difficulty determines the evidence, the evidence determines the posterior, and the posterior is the score. We derived where each difficulty tier is pedagogically useful and found that standard guess/slip parameters compress those points into the top of the mastery scale, leaving the hardest tier optimal only at complete mastery. We constructed a case in which no policy can change what is learned, and found the reported gain still separating by up to $19\%$ of its true value. And we found that the cost of mis-routing is asymmetric in a direction that depends on the learning model assumed, with the worst case of erring hard roughly three times that of erring easy, which makes conservative routing a minimax rather than a universally optimal choice.

None of these require abandoning simulation. They require measuring the learner rather than the belief about the learner, which in most existing harnesses is a matter of recording a variable that is already there.

\section*{Reproducibility}

All results are reproducible from seeded code released at \url{https://github.com/OmegaKavya/CogniLoop}. The simulation harness, the policy specifications, and the scripts producing each table are included, together with a test suite that verifies the knowledge-tracing implementation used in the simulation is numerically identical to the one in the deployed system it was derived from.

\bibliographystyle{acmtrans}
\bibliography{estimator_confound}

\end{document}
